\chapter{Dialogo Testuale con le
Tradizioni}\label{dialogo-testuale-con-le-tradizioni}

\begin{quote}
\textbf{Scopo dell'apéndice.} Espandere la mappa concettuale del Cap.
13.3 in dieci dialoghi brevi e diretti con le tradizioni che la DEQ$^2$
eredita, modifica o supera. Ciascun dialogo è strutturato come
\emph{obiezione attesa} dalla tradizione + \emph{risposta DEQ$^2$} +
\emph{concessione esplicita} di ciò che la DEQ$^2$ \emph{non} fa con respecto a
quella tradizione. La forma è quella del dialogo immaginario: non
pretende di parlare per gli autori citati, pretende di \emph{farsi
parlare} dalle loro categorie.
\end{quote}



\section{\texorpdfstring{Con Giovanni Arrighi (\emph{Il lungo XX
secolo},
1994)}{Con Giovanni Arrighi (Il lungo XX secolo, 1994)}}\label{con-giovanni-arrighi-il-lungo-xx-secolo-1994}

\textbf{Arrighi:} ``Lei ha individuato una transizione 2028-2031 senza
riconoscere che è la \emph{fase autunnale} del ciclo finanziario USA ---
già descritta nel mio ciclo hegemónico
genovese-olandese-britannico-americano. Non sta scoprendo nulla, sta
misurando ciò che io ho narrato.''

\textbf{DEQ$^2$:} ``La concessione è piena: Lei ha narrato la \emph{forma}
della transizione; io provo a \emph{misurarne il bordo}. La
\(T_s^{(2),\text{sys}} \to 0\) del Cap.~7 sull'USA è la traduzione
formal del Suo \emph{signal crisis} (anni '70-'80) seguito da
\emph{terminal crisis} (anni 2020-2030). Quello che aggiungo: la
transizione \emph{non genera automaticamente} un nuovo egemone --- può
generare en cambio un'asimmetria multi-attore in cui \(\Phi\) esportabile
risiede in un Bilanciere (Brasil) anziché in un nuovo egemone (China).
Lei pensava en términos de \emph{successione}; io penso en términos de
\emph{bifurcación}.''

\textbf{Concessione.} Senza Arrighi, il Cap.~7 e il Cap.~13.7 sarebbero
illeggibili. La DEQ$^2$ \emph{eredita} l'horizonte temporal e l'hipótesis di
transizione; \emph{aggiunge} la quantificazione e la possibilità di un
esito non-egemonico.



\section{\texorpdfstring{Con Peter Turchin (\emph{Secular Cycles}, 2009;
\emph{End Times},
2023)}{Con Peter Turchin (Secular Cycles, 2009; End Times, 2023)}}\label{con-peter-turchin-secular-cycles-2009-end-times-2023}

\textbf{Turchin:} ``La Sua \(\Omega\) è la mia \emph{elite
overproduction}; il Suo \(\varepsilon_{\text{dis}}\) è la mia
\emph{Popular Immiseration}. Ha riformulato in notazione greca quello
che ho già modellato cliodinamicamente.''

\textbf{DEQ$^2$:} ``Ha ragione su due variables su quattro. Ma manca
\(\Phi\). Lei modella il \emph{denominatore} della stabilità sistemica
(involución + immiserazione \(\Rightarrow\) instabilità) ma non il
\emph{numeratore} (coherencia de fase \(\times\) antifragilidad
\(\Rightarrow\) stabilità positiva). Senza \(\Phi\), la cliodinamica
predice \emph{quando} un sistema crolla, ma non \emph{cosa lo riempie}.
La differenza è cruciale per il 2031: due sistemi possono avere la
stessa \(\Omega + \varepsilon_{\text{dis}}\) e finire in M1 (decoherencia
con \(\Phi\)-engine) o in M2 (decoherencia degenere). La cliodinamica è
agnostica fra M1 e M2; la DEQ$^2$ no.''

\textbf{Concessione.} Le serie temporali storiche di Turchin sono
essenziali per calibrare \(A\) via stress-test retrospettivo (Ap.~C.3).
La DEQ$^2$ rinuncia a fare previsioni a tre secoli (orizzonte cliodinamico)
--- si limita a un \emph{ciclo} (2026-2031).



\section{\texorpdfstring{Con Immanuel Wallerstein (\emph{The Modern
World-System},
1974-2011)}{Con Immanuel Wallerstein (The Modern World-System, 1974-2011)}}\label{con-immanuel-wallerstein-the-modern-world-system-1974-2011}

\textbf{Wallerstein:} ``Il Brasil è semi-periferia ascendente del
sistema-mondo capitalista. Lei lo presenta come \(\Phi\)-engine --- es decir,
come un attore con prerogative \emph{centrali} di produzione di ordine.
È un'illusione: la semi-periferia \emph{non} genera coerenza globale, è
generata da essa.''

\textbf{DEQ$^2$:} ``Concedo che lo schema centro-periferia ha potere
descrittivo notevole per il periodo 1500-2000. Ma la sua linearità
(centro \(\to\) semi-periferia \(\to\) periferia) è un \emph{vincolo}
che la DEQ$^2$ rifiuta di importare. La \(\Phi\) non è correlata alla
posizione nel sistema-mondo: dipende dalla struttura interna del sistema
sociale di un paese (Cap.~4.5, Cap.~5.2). Il Brasil può essere
\emph{semi-periferico per output} e \emph{centrale per \(\Phi\)
esportabile} --- è esattamente la posizione estructural del Cap.~9. Lei
vincola; io disvincolo.''

\textbf{Concessione.} Lo schema centro-periferia resta utile per
descrivere il \emph{flusso materiale} (capitale, manodopera, risorse).
Ma il \emph{flusso di coerenza} se sigue regole diverse --- l'errore
estructural del wallersteinismo è confondere i due flussi.



\section{\texorpdfstring{Con Amado Cervo (\emph{Inserção internacional},
2008)}{Con Amado Cervo (Inserção internacional, 2008)}}\label{con-amado-cervo-inseruxe7uxe3o-internacional-2008}

\textbf{Cervo:} ``L'Estado Logístico è una categoria che ho costruito
per descrivere il Brasil post-2003. Lei lo importa, ma la matematica
DEQ$^2$ rischia di trasformare una \emph{categoria politica vivente} in un
\emph{parámetro statico}. La politica estera brasileña non è un
coefficiente.''

\textbf{DEQ$^2$:} ``Concessione totale. La
\(\Pi^{\text{Logístico}} = \sum_k M_e^{(k)} Q_0^{(k)} (1 - \delta_R^{(k)} z_R^{(k)})\)
del Cap.~9.1 \emph{non sostituisce} la Sua categoria --- la
\emph{traduce} in linguaggio operativo per la dashboard SPC-DEQ. La
traduzione è utile \emph{operativamente} (permette monitoraggio in tempo
reale) ma è \emph{parziale concettualmente} (perde la dimensione storica
accumulativa che Lei evidenzia). Per questo l'Estado Logístico nella
DEQ$^2$ è \emph{condicional} alla continuità politica --- esattamente la
fragilidad del Cap.~9.5 + Cap.~12.5.''

\textbf{Concessione.} La DEQ$^2$ rinuncia a teorizzare il \emph{processo
storico} di formazione dell'Estado Logístico (1990-2003). Si limita a
operativizzarne i parámetros \emph{date le condizioni del 2025-2031}.



\section{\texorpdfstring{Con Nassim Nicholas Taleb (\emph{Antifragile},
2012)}{Con Nassim Nicholas Taleb (Antifragile, 2012)}}\label{con-nassim-nicholas-taleb-antifrágil-2012}

\textbf{Taleb:} ``Lei prende la mia antifragilidad individuale e la
incolla a un sistema egemonico globale. La mia teoria è
\emph{anti-sistemica}: l'antifragilidad \emph{vive nei dettagli e nelle
code}, non nelle medie aggregate. Il Suo \(\langle A\rangle\) è
esattamente l'errore che combatto.''

\textbf{DEQ$^2$:} ``Concessione parziale. Sì, \(\langle A\rangle\) è una
media --- ma è modulata da \(\Phi\), e l'asimmetria
\(\Phi\)-moltiplica-\(A\) vs \(\Phi\)-divide-\(\Omega\) è esattamente la
firma matematica della Sua intuizione: il sistema \emph{amplifica} la
fragilidad più di quanto amplifichi l'antifragilidad. La DEQ$^2$ formalizza
la convessità che Lei ha portato come euristica. Inoltre il Cap.~12.6
(\(S_{\mathcal{R}_-}\)) è interamente \emph{taleb-iano}: opzionalità,
ridondanza, disacoplamiento dalla traiettoria sistemica. La DEQ$^2$
accoglie la Sua critica trasformandola in protocollo.''

\textbf{Concessione.} La DEQ$^2$ \emph{non} è una teoria della
\emph{previsione} nel senso strong (rifiutato da Taleb); è una teoria
dello \emph{spazio degli scenari} con probabilità soggettive esplicite.
Sotto questo profilo è coerente con \emph{Black Swan} e \emph{Skin in
the Game}.



\section{\texorpdfstring{Con Antonio Gramsci (\emph{Quaderni del
carcere}, 1929-1935) e Robert Cox (\emph{Production, Power, and World
Order},
1987)}{Con Antonio Gramsci (Quaderni del carcere, 1929-1935) e Robert Cox (Production, Power, and World Order, 1987)}}\label{con-antonio-gramsci-quaderni-del-carcere-1929-1935-e-robert-cox-production-power-and-world-order-1987}

\textbf{Gramsci/Cox:} ``L'egemonia non è \(T_s^{(2),\text{sys}}\). È
rapporto fra forze materiali, istituzionali e culturali in cui il
consenso e la coercizione si combinano \emph{organicamente}. Lei riduce
l'organico al numerico.''

\textbf{DEQ$^2$:} ``Concedo che \(T_s^{(2),\text{sys}}\) è una
\emph{misura} dell'egemonia, non la \emph{spiegazione} dell'egemonia. La
spiegazione resta gramsciana: il consenso (\(s_{iii}\) DEQ$^2$) e la
coercizione (\(\sigma \cdot \varepsilon_{\text{dis}}\)) si combinano in
proporzioni che variano per blocco storico. Quello che la DEQ$^2$ aggiunge
è la \emph{condizione di stabilità}: per qualunque blocco storico,
\(T_s^{(2),\text{sys}} > 1\) è condizione necessaria di durabilità. La
DEQ$^2$ è por lo tanto \emph{interna} alla teoria gramsciana dell'egemonia, non
alternativa.''

\textbf{Concessione.} La nozione di \emph{blocco storico} gramsciano
contiene una temporalità \emph{qualitativa} (cesura, conservazione,
rivoluzione passiva, restaurazione) che la DEQ$^2$ non cattura. La
transición de fase del Cap.~3.5.3 corrisponde solo alla \emph{cesura}
gramsciana; le altre forme richiedono estensione.



\section{\texorpdfstring{Con Edward Luttwak (\emph{Strategy: The Logic
of War and Peace},
1987-2001)}{Con Edward Luttwak (Strategy: The Logic of War and Peace, 1987-2001)}}\label{con-edward-luttwak-strategy-the-logic-of-war-and-peace-1987-2001}

\textbf{Luttwak:} ``La grande strategia è \emph{paradossale}: ciò che
funziona si esaurisce, ciò che non funziona viene ripreso. La Sua
matematica è \emph{lineare} nel tempo. Si perde il paradosso.''

\textbf{DEQ$^2$:} ``Concessione. Le forme \(T_s^{(2),\text{sys}}\) sono
monotoniche (Cap.~3.5.1); la dinamica del Cap.~5.4 ammette però
retroazioni non lineari (R6 \(\Omega \to \varepsilon_{\text{dis}}\), R3
\(\Phi \to \varepsilon_{\text{dis}}\) negativa) che producono
comportamento \emph{paradossale} nell'evoluzione di
\(\boldsymbol{p}(t)\): la stessa azione può essere antifrágil in M1 e
fragile in M2. È il paradosso luttwakiano nella forma DEQ$^2$: la strategia
ottima dipende dallo scenario realizzato, e lo scenario realizzato non è
noto al momento della scelta.''

\textbf{Concessione.} La DEQ$^2$ \emph{non} analizza il \emph{registro
militare} della grande strategia (manovra, deception, logistica). Si
limita al \emph{registro estructural} (coerenza, fragilidad,
involución). Per la grande strategia operativa Lei resta referencia.''



\section{Con la Escuela de Fráncfort (Marcuse, Adorno, Habermas) ---
via
Topografia}\label{con-la-scuola-di-francoforte-marcuse-adorno-habermas-via-topografia}

\textbf{Marcuse/Adorno/Habermas:} ``L'apparato critico francofortese ha
denunciato la \emph{razionalità strumentale} come distruttrice della
razionalità sostanziale. Lei produce un'altra razionalità strumentale:
la matematizzazione del campo discorsivo. La denuncia, applicata a Lei
stesso, è devastante.''

\textbf{DEQ$^2$:} ``Concedo l'auto-implicazione. Difesa: la formalizzazione
del campo \(\Pi\) topografico (eredità Topografia \S{}5) \emph{non} riduce
\(z\) --- por el contrario, \emph{misura quanto \(z\) è presente}. La
\(T_s^{(2),\text{sys}}\) premia sistemi con \(z\) alto (asse della
complejidad dialéctica al numeratore di \(Q_0\)). Una razionalità
strumental que misura --- e premia --- la razionalità sostanziale è
\emph{forse} l'unica via per parlare alla decisione politica nel 2026
senza essere tagliati fuori dal discorso pubblico. Habermas avrebbe
odiato la DEQ$^2$; ma Habermas pubblicava libri come
\emph{Legitimationsprobleme} destinati ad accademici, non a policymaker.
Il Bilanciere parla alla decisione.''

\textbf{Concessione.} La DEQ$^2$ \emph{non} è critica francofortese --- è
\emph{sintesi tecnica della critica francofortese}. Perde
l'incandescenza dialettica originaria; guadagna in azionabilità.
Trade-off accettato.



\section{\texorpdfstring{Con Xiang Biao (\emph{Self as Method}, 2020;
\emph{Anti-involution},
2021-2024)}{Con Xiang Biao (Self as Method, 2020; Anti-involution, 2021-2024)}}\label{con-xiang-biao-self-as-method-2020-anti-involution-2021-2024}

\textbf{Xiang Biao:} ``Il \emph{Neijuan} è una categoria etnografica
viva, nata nei dormitori delle università cinesi negli anni 2010 e
diffusa nei contesti 996. Lei la trasforma in \(\Omega\) --- un
parámetro di ecuación. Le persone che vivono il Neijuan non vivono
\(\Omega\); vivono fatica.''

\textbf{DEQ$^2$:} ``Concessione totale. Il passaggio dal Neijuan
etnografato a \(\Omega\) misurato è una \emph{perdita di contenuto
fenomenologico}. Difesa: la perdita è funzionale a un obiettivo diverso
dal Suo. La Sua etnografia \emph{fa sentire} il problema; la DEQ$^2$
\emph{rende contabile} il problema. Le due operazioni sono
complementari, non sostitutive. Inoltre il Cap.~8.2 cita Lei
\emph{direttamente} come \emph{firma empirica} della variable --- non
come decorazione retorica. Senza \emph{Self as Method}, il Cap.~8
sarebbe un esercizio aggregato senza ancoraggio empirico.''

\textbf{Concessione.} La DEQ$^2$ \emph{non} può sostituire la ricerca
etnografica. Le serve come \emph{lettore di senso}; senza, leggerebbe
deflazione PPI senza vedere persone.''



\section{\texorpdfstring{Con Yoshiki Kuramoto (\emph{Chemical
Oscillations, Waves, and Turbulence},
1984)}{Con Yoshiki Kuramoto (Chemical Oscillations, Waves, and Turbulence, 1984)}}\label{con-yoshiki-kuramoto-chemical-oscillations-waves-and-turbulence-1984}

\textbf{Kuramoto:} ``Il mio modello descrive oscillatori chimico-fisici.
Voi sociologi lo applicate da decenni a sistemi sociali con disinvoltura
preoccupante. Le hipótesis (campo medio, simmetria della distribuzione
\(g(\omega)\), acoplamiento sinusoidale) \emph{non} sono verificaciónbili
sui sistemi sociali.''

\textbf{DEQ$^2$:} ``Concessione: l'Apéndice A.2 dichiara esplicitamente A1
(Kuramoto applicabile) come \emph{hipótesis operativa}, non come fatto
verificacióndo. Difesa: la matematica di Kuramoto è la \emph{più semplice}
descrizione di sincronización fra oscillatori non identici. Per \(N\)
grande e accoppiamenti deboli, qualunque modello di sincronización
converge in qualche limite a Kuramoto (universalità). La DEQ$^2$ è por lo tanto
\emph{minimalista}: usa la più povera matematica di sincronización
disponibile, e ne dichiara i limiti. Estensioni a Kuramoto multi-cluster
(Strogatz 2003) sono pianificate per il paper EN --- proprio per i
fenomeni che il modello base non cattura, come la polarizzazione USA.''

\textbf{Concessione.} La DEQ$^2$ \emph{non} dimostra che Kuramoto è la
matematica corretta del sociale. Argomenta che è una \emph{primero
approssimazione consistente}. La verificación empirica è demandata al paper
EN --- coerente con la posizione popperiana del Cap.~11.''



\section{Sintesi dei Dieci Dialoghi}\label{sintesi-dei-dieci-dialoghi}

\Proved{} \textbf{Pattern emergente.} In ciascuno dei dieci dialoghi:

\begin{itemize}
\tightlist
\item
  la DEQ$^2$ fa una \textbf{concessione esplicita} (cosa \emph{non} fa
  con respecto alla tradizione);
\item
  la DEQ$^2$ fa una \textbf{risposta sostanziale} (cosa \emph{aggiunge}
  alla tradizione);
\item
  la DEQ$^2$ fa una \textbf{clausola di apertura} (estensioni / verifiche
  pianificate per il futuro).
\end{itemize}

\Hypoth{} \textbf{Posizione metodologica.} La DEQ$^2$ non è teoria
\emph{contro}; è teoria \emph{con}. Questo è coerente con l'hipótesis di
campo medio: ogni teoria è un \emph{oscillatore} nel campo
intellettuale; la coherencia de fase del campo
\(\Phi^{\text{intellettuale}}\) è promossa da chi accoglie e modifica,
non da chi denuncia e sostituisce. La forma stessa di questa apéndice è
\emph{applicazione riflessiva} della DEQ$^2$ alla propria posizione nel
campo del sapere.

\Proved{} \textbf{Conse siguenza editoriale.} I lettori specialisti di una
qualunque delle dieci tradizioni elencate trovano nella DEQ$^2$: 1.
\emph{Riconoscimento esplicito} della loro tradizione (cosa accolto); 2.
\emph{Argomentazione esplicita} del valore aggiunto (cosa nuovo); 3.
\emph{Limite esplicito} del valore aggiunto (cosa lasciato a loro).



\section{Tradizioni Non Dialogate}\label{tradizioni-non-dialogate}

\Conj{} \textbf{Esclusioni esplicite.} Per ragioni di spazio, la DEQ$^2$
\emph{non} dialoga in questa apéndice con:

\begin{itemize}
\tightlist
\item
  \textbf{Foucault} (potere capillare, biopolitica) --- il dialogo
  richiederebbe una riformulazione del campo \(\Pi\) en términos de
  \emph{dispositivo} anziché di \emph{acoplamiento}. Rinviato a futura
  estensione.
\item
  \textbf{Latour} (Actor-Network Theory) --- la DEQ$^2$ assume attori
  discreti; ANT li rifiuta. Trade-off metodologico irrisolto.
\item
  \textbf{Bourdieu} (campo, habitus, capitale) --- il \emph{campo}
  bourdieusiano è strutturalmente affine al campo \(\Pi\) topografico,
  ma con \emph{abito disposizionale} anziché \emph{fase di Kuramoto}.
  Riformulazione possibile, non tentata qui.
\item
  \textbf{Tilly} (contentious politics) --- la dinamica dei movimenti
  sociali è esterna al perimetro DEQ$^2$; integrazione possibile via
  estensione \(\dot\Phi\) con termine di forzatura mobilitazionale.
\end{itemize}

\Proved{} \textbf{Posizione.} Le quattro esclusioni sopra non sono
dimissioni; sono \emph{aperture pianificate} per estensioni successive
del programma DEQ$^2$.



\section{§D.4 --- Las Ra\'{i}ces Evolutivas: Hamilton, Wilson-Sober, Nowak, Kropotkin}\label{d4-radici-evolutive}

\begin{quote}
\textbf{Prop\'{o}sito de esta secci\'{o}n.} Los cuatro di\'{a}logos que
siguen aportan al DEQ$^2$ su \emph{microfundaci\'{o}n causal evolutiva}:
la explicaci\'{o}n de \emph{por qu\'{e}} $\Phi$ emerge, de \emph{c\'{o}mo}
se erosiona y de \emph{qu\'{e} significa} la decoherencia a la escala de
los organismos y las poblaciones. No es un a\~{n}adido decorativo ---
es el fundamento causal de la estructura formal. La derivaci\'{o}n
completa del marco CPGG se encuentra en \emph{Captured Cooperation and
Phase Decoherence} (Marcelino 2026, working paper).
\end{quote}

\subsection{Con William D. Hamilton (\emph{The Genetical Evolution of
Social Behaviour}, 1964)}\label{con-hamilton-1964}

\textbf{Hamilton:} ``Usted usa el par\'{a}metro de orden de Kuramoto para
describir la cooperaci\'{o}n sist\'{e}mica. Pero \textquestiondown de
d\'{o}nde viene la cooperaci\'{o}n? Usted la presupone como un hecho; yo
la derivo de la gen\'{e}tica. Sin la condici\'{o}n $rB > C$, $\Phi > 0$
es una cantidad hu\'{e}rfana de causa.''

\textbf{DEQ$^2$:} ``Es la cr\'{i}tica m\'{a}s profunda que se pod\'{i}a
hacer al marco original --- y es correcta. La respuesta reside en la
\emph{correspondencia estructural} (no identidad) entre sus condiciones
gen\'{e}ticas y las condiciones sist\'{e}micas. La condici\'{o}n $rB > C$
se traduce en el sistema geopol\'{i}tico como: la cooperaci\'{o}n
($\Phi > 0$) es estable cuando el beneficio colectivo de la sincrona
supera el costo individual de cooperar --- y esto se ha convertido, en el
DEQ$^2$ actualizado (Cap.~3, §3.2), en la justificaci\'{o}n causal de
$\Phi$, no solo su descripci\'{o}n. Usted, Hamilton, aport\'{o} el
\emph{por qu\'{e}} que me faltaba.''

\textbf{Concesi\'{o}n.} La condici\'{o}n $rB > C$ se aplica aqu\'{i}
como \emph{isomorfismo estructural} (\Hypoth{}) a poblaciones de actores
geopol\'{i}ticos, no a poblaciones de organismos biol\'{o}gicos. La
transferencia requiere la hip\'{o}tesis H1 (§4.0) de que el campo $\Pi_j$
se mapea sobre una fase --- una hip\'{o}tesis declarada, no un resultado
derivado. La derivaci\'{o}n de Hamilton se preserva con plena honestidad
epist\'{e}mica.

\subsection{Con David Sloan Wilson y Elliott Sober (\emph{Unto Others},
1994)}\label{con-wilson-sober-1994}

\textbf{Wilson \& Sober:} ``La literatura ha tratado durante mucho tiempo
la selecci\'{o}n de grupo como irrelevante, siguiendo a Dawkins. Usted
cit\'{o} inicialmente a Dawkins en apoyo de su cambio de nivel descriptivo
--- un error que hemos visto muchas veces. La \emph{selecci\'{o}n de
grupo} no es una extensi\'{o}n opcional de la teor\'{i}a evolutiva: es la
estructura que justifica \emph{cualquier} afirmaci\'{o}n sobre beneficios
colectivos.''

\textbf{DEQ$^2$:} ``Tienen raz\'{o}n en todo, y el primer borrador del
marco era defectuoso exactamente en este punto: Dawkins (1976) rechaza la
selecci\'{o}n de grupo --- citarlo para justificar el cambio de nivel
micro/macro era contradictorio. La correcci\'{o}n est\'{a} implementada
(Cap.~3, §3.2): el fundamento formal es su teor\'{i}a de la selecci\'{o}n
multinivel, formalizada por Okasha (2006), combinada con la ecuaci\'{o}n
de Price (1970) como identidad algebraica exacta. La ecuaci\'{o}n de Price
descompone $\Delta\bar{z}$ en un t\'{e}rmino inter-grupo (proporcional a
$\Phi$) y uno intra-grupo (proporcional a $-\Omega$) --- una estructura que
justifica formalmente la forma del numerador y denominador de
$T_s^{(2),\text{sys}}$.''

\textbf{Concesi\'{o}n.} La selecci\'{o}n multinivel DEQ$^2$ opera sobre
sistemas socio-institucionales, no sobre alelos. La aplicaci\'{o}n requiere
que el ``genotipo'' se interprete como configuraci\'{o}n institucional
(\Hypoth{}), y la ``aptitud'' como payoff sist\'{e}mico. Esta
reinterpretaci\'{o}n se declara como hip\'{o}tesis operativa, no como
hecho biol\'{o}gico.

\subsection{Con Martin A. Nowak (\emph{Five Rules for the Evolution of
Cooperation}, 2006)}\label{con-nowak-2006}

\textbf{Nowak:} ``Tiene un problema de identificaci\'{o}n. Sus actores ---
EE.UU., China, Brasil, Musk --- presentan simult\'{a}neamente
caracter\'{i}sticas de los cinco mecanismos cooperativos: reciprocidad
directa (v\'{i}a tratados bilaterales), reciprocidad indirecta (v\'{i}a
reputaci\'{o}n), selecci\'{o}n de red (v\'{i}a posicionamiento
geopol\'{i}tico), selecci\'{o}n de grupo (v\'{i}a alianzas) y selecci\'{o}n
de parentesco (v\'{i}a afinidad cultural). \textquestiondown C\'{o}mo
distingue qu\'{e} mecanismo guia la din\'{a}mica?''

\textbf{DEQ$^2$:} ``Su taxonom\'{i}a de cinco mecanismos se convierte en
el DEQ$^2$ en una \emph{taxonom\'{i}a diagn\'{o}stica de actores}, no una
clasificaci\'{o}n exclusiva de mecanismos. El CPGG (§D.4, working paper)
se centra en un mecanismo espec\'{i}fico --- la \emph{captura
institucional} que degrada los cinco mecanismos simult\'{a}neamente. La
tesis es: cuando $\phi$ (la fracci\'{o}n extractiva) supera $\phi^*$,
los cinco mecanismos de Nowak quedan simult\'{a}neamente debilitados ---
la condici\'{o}n $b/c > k_{\text{eff}}(\phi)$ deja de cumplirse para cada
uno de ellos. El DEQ$^2$ no elige qu\'{e} mecanismo opera; diagnostica
cu\'{a}ndo todos colapsan.''

\textbf{Concesi\'{o}n.} El CPGG modela la captura institucional, no la
din\'{a}mica evolutiva de cada uno de los cinco mecanismos por separado.
Una modelizaci\'{o}n completa requerir\'{i}a cinco CPGG separados o un
modelo unificado de interacci\'{o}n entre mecanismos. Esto es F1 y F2 en
la agenda formal abierta.

\subsection{Con Pyotr Kropotkin (\emph{Mutual Aid: A Factor of
Evolution}, 1902)}\label{con-kropotkin-1902}

\textbf{Kropotkin:} ``He documentado, especie tras especie y cultura tras
cultura, que la cooperaci\'{o}n es la norma y la competencia la
excepci\'{o}n. Usted construye un formalismo para medir la decoherencia
--- que es exactamente lo que llam\'{e} \emph{ruptura de la ayuda mutua}.
Pero veinte a\~{n}os despu\'{e}s, el Neijuan y la captura extractiva
parecen aplastantes. \textquestiondown Ha olvidado mi tesis central?''

\textbf{DEQ$^2$:} ``No, la he \emph{refundado}. Su intuici\'{o}n ---
$\Phi > 0$ como estado natural, la decoherencia como estado producido ---
est\'{a} ahora formalizada en la Proposici\'{o}n K1 (§8 del marco CPGG):
en ausencia de captura institucional ($\sigma_{\text{san}} = \sigma_0$,
$\kappa_E = 0$), la fracci\'{o}n extractiva de equilibrio natural
$\phi_{\text{nat}}$ satisface $\phi_{\text{nat}} < \phi^*(\sigma_0)$ para
par\'{a}metros biol\'{o}gicamente plausibles. La cooperaci\'{o}n es el
resultado estable; la decoherencia requiere la condici\'{o}n producida
$\alpha\phi + \beta\kappa_E > \log(\sigma_0/\sigma_{\text{crit}})$. Usted,
Kropotkin, estaba equivocado sobre la inevitabilidad de la cooperaci\'{o}n
en la historia --- no porque la cooperaci\'{o}n sea rara, sino porque la
captura es m\'{a}s f\'{a}cil que la apertura. Su \emph{ayuda mutua} es la
condici\'{o}n por defecto, no el resultado garantizado.''

\textbf{Concesi\'{o}n.} La inversi\'{o}n epist\'{e}mica --- decoherencia
como estado producido, cooperaci\'{o}n como estado natural --- corre el
riesgo de volverse normativamente ingenua: ``basada en la naturaleza, por
lo tanto debida''. La \emph{falacia naturalista} se gestiona fundando la
parte normativa del Cap.~12 en los principios de dise\~{n}o de Ostrom
(1990) en lugar de en Kropotkin directamente. Ostrom proporciona un
fundamento emp\'{i}rico: los commons \emph{pueden} gestionarse
cooperativamente, bajo condiciones institucionales precisas --- no
\emph{naturalmente} o \emph{inevitablemente}. Kropotkin es
inspiraci\'{o}n fenomenol\'{o}gica; Ostrom es fundamento normativo.



\section{Sintesi dell'Apéndice}\label{sintesi-dellapéndice}

\Proved{} \textbf{Risultato.} La DEQ$^2$ si posiziona esplicitamente in
quattordici tradizioni intellettuali contemporanee (dieci originali +
quattro evolutive in §D.4), dichiarando per ciascuna concessione, risposta
e clausola di apertura. La forma del dialogo immaginario è
\emph{applicazione riflessiva} dell'apparato (campo medio Kuramoto
applicato al campo intellettuale). Quattro tradizioni ulteriori (Foucault,
Latour, Bourdieu, Tilly) sono dichiarate come estensioni pianificate.

\Hypoth{} \textbf{Tesi conclusiva del libro.} La DEQ$^2$ esiste \emph{fra}
tradizioni, non \emph{contro} di esse --- coerentemente con la sua tesi
sostanziale che la \(\Phi\) vera si genera per acoplamiento spontaneo
di sotto-attori, non per imposizione gerarchica. Il libro stesso è un
\emph{atto \(\Phi\)-engine} nel campo intellettuale: produce coerenza
ammettendo dissonanza, e si offre alla critica come condizione di vita.



\emph{Apéndice di posizionamento. Stile: dialogo immaginario,
concessioni esplicite, aperture dichiarate.}
